\section{Limitations}
\label{sec:limitations}

The experiments explore a limited part of the approach's potential. We train the denoisers from
scratch, without the text pretraining used by several competing speech--text systems. This may
particularly affect speech recognition and the semantic coherence of continuations, where paired
speech data must supply linguistic knowledge as well as cross-modal alignment. Text-pretrained
initialisation is therefore a promising direction, although controlled comparisons are needed to
establish its contribution. Protein-sequence pretraining may offer analogous benefits for
sequence--structure modelling.

Our models are also generally smaller than the strongest systems for the corresponding tasks,
and the training datasets provide limited coverage of the distributions we aim to model. For
example, the speech experiments rely on read speech, offering relatively narrow linguistic and
acoustic diversity. These choices restrict what the present results can establish about
performance at scale. Larger backbones and broader, higher-quality datasets are plausible routes
to improvement, but their effects need to be measured within this formulation.

Continuous modalities introduce a further dependence on pretrained tokenizers and decoders.
Information discarded during tokenisation is unavailable to the denoiser, while reconstruction
errors can persist even when codes are modelled accurately. The observed performance therefore
reflects both the bitstream model and its representation. Codec reconstruction results help
distinguish these effects, but do not by themselves explain the remaining generation errors.

\section{Discussion and conclusion}
\label{sec:discussion}
\label{sec:conclusion}

This work demonstrates that a flattened one-dimensional bitstream can serve as a common
representation for multimodal generation across three distinct forms of cross-modal
dependence. On MNIST-Sum, raw pixel bits and symbolic text bits form one learnable
joint distribution with no learned tokenizer: 97.56\% exact image$\rightarrow$text, 99.95\%
text$\rightarrow$image and 99.46\% mutually consistent joint samples. On speech and text, the same
formulation handles a strongly asymmetric, codec-backed pair, co-generating speech and text at
3.5\% cross-modal WER and UTMOS 3.77 from a model trained without text pretraining, while speech
recognition and thematic continuation remain its weakest directions. On protein sequence and
structure, one denoiser performs folding, inverse folding and joint generation of physically
coupled channels, and fold-level self-consistency is well ahead of coordinate-level precision:
68.75\% (55/80) of generated backbones reach scTM~$\geq$~0.9, against 38.75\% (31/80) at
scRMSD~$\leq$~2\,\AA. With changes largely
confined to tokenisation and layout, one continuous denoising formulation thus carries across a
controlled pair, an asymmetric pair and a physically coupled pair; within each checkpoint, joint
and conditional generation share the same weights.

The breadth of these results, obtained with modest backbones trained from scratch, supports the
feasibility of the approach. It motivates investigating how far stronger pretraining, larger
models, improved data coverage and better continuous-modality tokenizers can advance performance
while preserving the shared bitstream formulation.
